% GENERATED FILE. Edit canonical lessons, then run npm run generate:books.
% canonical-source-hash: fnv1a64:b3bfbdb1202ac44e
% canonical-lessons: GU-C237-kagdo, GU-C237-ghuvad, GU-C237-pakh, GU-C237-phaphadvu, GU-C237-udvu

\chapter{Wings in the air}
\label{ch:gu-wings-in-the-air}
\hlchaptermodality{237}

\begin{chapteropening}
\textbf{By the end of this chapter:} \emph{I can say how a bird flew away.}

Complete the last lesson of chapter 237: I can say how a bird flew away.
\end{chapteropening}

\section[\texorpdfstring{\gu{કાગડો} (kāgḍo)}{કાગડો (kāgḍo)}]{\gu{કાગડો} (kāgḍo) — crow}

\label{lesson:GU-C237-kagdo}



\begin{quote}
Before the new one: say the Gujarati for beak, then the Gujarati for to peck.
\end{quote}

\subsection*{You'll want to know: \gu{કાગડો}}
\textbf{\gu{કાગડો}} — \emph{kāgḍo} — \textquotedblleft{}crow\textquotedblright{}.

\begin{quote}
\textbf{\gu{ગઈકાલે} \gu{એક} \gu{કાગડો} \gu{ઘરમાં} \gu{આવ્યો}.} — \emph{gaīkāle ek kāgḍo gharmā̃ āvyo} — \textquotedblleft{}Yesterday a crow came into the house.\textquotedblright{}
\end{quote}

\begin{cousinweb}
\textbf{\gu{કાગડો}} (\emph{kāgḍo}) is usually traced to Sanskrit \emph{kāka}, \textquotedblleft{}crow\textquotedblright{}, an imitative word that CDIAL says took various suffixes, with \emph{-k-} or \emph{-g-}, in the modern languages. CDIAL lists only Gujarati \emph{kāu}, \textquotedblleft{}crow\textquotedblright{}, there, so the path to \emph{kāgḍo} is not spelled out (CDIAL 2993).
\end{cousinweb}

\begin{grammarlens}[title={What you've built}]
One more word for everyday life at A2.
\end{grammarlens}

\subsection*{Your turn}
\begin{itemize}
\raggedright
  \item \emph{Say it:} \emph{kāgḍo}
  \item \emph{Say it:} \emph{kāgḍo}, once more
  \item \emph{Say it:} \emph{gaīkāle ek kāgḍo gharmā̃ āvyo} — Yesterday a crow came into the house.
  \item \emph{Recall:} say the Gujarati for beak, then the Gujarati for to peck, then say \emph{kāgḍo} again
\end{itemize}

\begin{tcolorbox}[breakable,colback=teal!4,colframe=teal!35!black,title={Before you move on}]
What does \gu{કાગડો} mean? (\textquotedblleft{}Crow\textquotedblright{}.) Say it once more.
\end{tcolorbox}
\section[\texorpdfstring{\gu{ઘુવડ} (ghuvaḍ)}{ઘુવડ (ghuvaḍ)}]{\gu{ઘુવડ} (ghuvaḍ) — owl}

\label{lesson:GU-C237-ghuvad}



\begin{quote}
Before the new one: say the Gujarati for to peck, then the Gujarati for crow.
\end{quote}

\subsection*{You'll want to know: \gu{ઘુવડ}}
\textbf{\gu{ઘુવડ}} — \emph{ghuvaḍ} — \textquotedblleft{}owl\textquotedblright{}.

\begin{quote}
\textbf{\gu{રાત્રે} \gu{ઘુવડ} \gu{બોલ્યું}.} — \emph{rātre ghuvaḍ bolyũ} — \textquotedblleft{}The owl called at night.\textquotedblright{}
\end{quote}

\begin{cousinweb}
\textbf{\gu{ઘુવડ}} (\emph{ghuvaḍ}) is inherited from Sanskrit \emph{ghūka}, \textquotedblleft{}owl\textquotedblright{}, which CDIAL takes from Dravidian, as in Tamil \emph{kūkai}, after T. Burrow (CDIAL 4494).
\end{cousinweb}

\begin{grammarlens}[title={What you've built}]
One more word for everyday life at A2.
\end{grammarlens}

\subsection*{Your turn}
\begin{itemize}
\raggedright
  \item \emph{Say it:} \emph{ghuvaḍ}
  \item \emph{Say it:} \emph{ghuvaḍ}, once more
  \item \emph{Say it:} \emph{rātre ghuvaḍ bolyũ} — The owl called at night.
  \item \emph{Recall:} say the Gujarati for to peck, then the Gujarati for crow, then say \emph{ghuvaḍ} again
\end{itemize}

\begin{tcolorbox}[breakable,colback=teal!4,colframe=teal!35!black,title={Before you move on}]
What does \gu{ઘુવડ} mean? (\textquotedblleft{}Owl\textquotedblright{}.) Say it once more.
\end{tcolorbox}
\section[\texorpdfstring{\gu{પાંખ} (pā̃kh)}{પાંખ (pā̃kh)}]{\gu{પાંખ} (pā̃kh) — wing}

\label{lesson:GU-C237-pakh}



\begin{quote}
Before the new one: say the Gujarati for crow, then the Gujarati for owl.
\end{quote}

\subsection*{You'll want to know: \gu{પાંખ}}
\textbf{\gu{પાંખ}} — \emph{pā̃kh} — \textquotedblleft{}wing\textquotedblright{}.

\begin{quote}
\textbf{\gu{પંખીની} \gu{પાંખ} \gu{ભીની} \gu{હતી}.} — \emph{pankhīnī pā̃kh bhīnī hatī} — \textquotedblleft{}The bird's wing was wet.\textquotedblright{}
\end{quote}

\begin{cousinweb}
\textbf{\gu{પાંખ}} (\emph{pā̃kh}) is a tadbhava, a word inherited from Sanskrit and reshaped by sound change. It comes from \emph{pakṣa}, \textquotedblleft{}wing, feather; side\textquotedblright{}; CDIAL queries where its nasal came from (CDIAL 7627).
\end{cousinweb}

\begin{grammarlens}[title={What you've built}]
One more word for everyday life at A2.
\end{grammarlens}

\subsection*{Your turn}
\begin{itemize}
\raggedright
  \item \emph{Say it:} \emph{pā̃kh}
  \item \emph{Say it:} \emph{pā̃kh}, once more
  \item \emph{Say it:} \emph{pankhīnī pā̃kh bhīnī hatī} — The bird's wing was wet.
  \item \emph{Recall:} say the Gujarati for crow, then the Gujarati for owl, then say \emph{pā̃kh} again
\end{itemize}

\begin{tcolorbox}[breakable,colback=teal!4,colframe=teal!35!black,title={Before you move on}]
What does \gu{પાંખ} mean? (\textquotedblleft{}Wing\textquotedblright{}.) Say it once more.
\end{tcolorbox}
\section[\texorpdfstring{\gu{ફફડવું} (phaphaḍvũ)}{ફફડવું (phaphaḍvũ)}]{\gu{ફફડવું} (phaphaḍvũ) — to flutter, to flap}

\label{lesson:GU-C237-phaphadvu}



\begin{quote}
Before the new one: say the Gujarati for owl, then the Gujarati for wing.
\end{quote}

\subsection*{You'll want to know: \gu{ફફડવું}}
\textbf{\gu{ફફડવું}} — \emph{phaphaḍvũ} — \textquotedblleft{}to flutter, to flap\textquotedblright{}.

\begin{quote}
\textbf{\gu{પંખીની} \gu{પાંખ} \gu{ફફડી}.} — \emph{pankhīnī pā̃kh phaphaḍī} — \textquotedblleft{}The bird's wing fluttered.\textquotedblright{}
\end{quote}

\begin{cousinweb}
\textbf{\gu{ફફડવું}} (\emph{phaphaḍvũ}) is imitative of the sound of wings. CDIAL lists the close Gujarati \emph{phaṛphaṛvũ}, \textquotedblleft{}to flutter\textquotedblright{}, under \emph{*phaṭ}, \textquotedblleft{}sudden movement\textquotedblright{}, a reconstructed form (CDIAL 9038).
\end{cousinweb}

\begin{grammarlens}[title={What you've built}]
One more word for everyday life at A2.
\end{grammarlens}

\subsection*{Your turn}
\begin{itemize}
\raggedright
  \item \emph{Say it:} \emph{phaphaḍvũ}
  \item \emph{Say it:} \emph{phaphaḍvũ}, once more
  \item \emph{Say it:} \emph{pankhīnī pā̃kh phaphaḍī} — The bird's wing fluttered.
  \item \emph{Recall:} say the Gujarati for owl, then the Gujarati for wing, then say \emph{phaphaḍvũ} again
\end{itemize}

\begin{tcolorbox}[breakable,colback=teal!4,colframe=teal!35!black,title={Before you move on}]
What does \gu{ફફડવું} mean? (\textquotedblleft{}To flutter, to flap\textquotedblright{}.) Say it once more.
\end{tcolorbox}
\section[\texorpdfstring{\gu{ઊડવું} (ūḍvũ)}{ઊડવું (ūḍvũ)}]{\gu{ઊડવું} (ūḍvũ) — to fly up, to fly}

\label{lesson:GU-C237-udvu}



\begin{quote}
Before the new one: say the Gujarati for wing, then the Gujarati for to flutter.
\end{quote}

\subsection*{You'll want to know: \gu{ઊડવું}}
\textbf{\gu{ઊડવું}} — \emph{ūḍvũ} — \textquotedblleft{}to fly up, to fly\textquotedblright{}.

\begin{quote}
\textbf{\gu{કાગડો} \gu{આકાશમાં} \gu{ઊડ્યો}.} — \emph{kāgḍo ākāśmā̃ ūḍyo} — \textquotedblleft{}The crow flew up into the sky.\textquotedblright{}
\end{quote}

\begin{cousinweb}
\textbf{\gu{ઊડવું}} (\emph{ūḍvũ}) is a tadbhava, a word inherited from Sanskrit and reshaped by sound change. It comes from \emph{uḍḍīyate}, \textquotedblleft{}flies up\textquotedblright{} (CDIAL 1697).
\end{cousinweb}

\begin{grammarlens}[title={What you've built}]
One more word for everyday life at A2.
\end{grammarlens}

\subsection*{Your turn}
\begin{itemize}
\raggedright
  \item \emph{Say it:} \emph{ūḍvũ}
  \item \emph{Say it:} \emph{ūḍvũ}, once more
  \item \emph{Say it:} \emph{kāgḍo ākāśmā̃ ūḍyo} — The crow flew up into the sky.
  \item \emph{Recall:} say the Gujarati for wing, then the Gujarati for to flutter, then say \emph{ūḍvũ} again
\end{itemize}

\begin{tcolorbox}[breakable,colback=teal!4,colframe=teal!35!black,title={Before you move on}]
What does \gu{ઊડવું} mean? (\textquotedblleft{}To fly up, to fly\textquotedblright{}.) Say it once more.
\end{tcolorbox}
